% !TEX root = ../main.tex
\chapter{Requisiti e casi d'uso}
\label{cap:requisiti}
In questo capitolo sono descritti gli attori del sistema, il confine di esso stesso e i casi d'uso principali, con l'obiettivo di chiarire le funzionalità del sistema e le interazioni con gli utenti.

\section{Attori e confine del sistema}
Gli attori del sistema sono gli utenti che interagiscono con esso. Gli attori principali sono:
\begin{itemize}
    \item \textbf{Visitor}: un utente non registrato che può registrare un nuovo account fornendo i dati anagrafici necessari e credenziali 
    valide oppure autenticarsi a un account esistente.
    \item \textbf{Utente registrato}: un utente non amministratore che ha creato un account nel sistema e può:
    \begin{itemize}
        \item Consultare venue, risorse e eventi;
        \item Modificare o eliminare il proprio account dopo la verifica della password;
        \item Cambiare la propria password;
        \item Presentare e gestire richieste di evento;
        \item Prenotare biglietti per eventi pubblici e pubblicati;
        \item Gestire le proprie prenotazioni, visualizzando lo stato della prenotazione e la possibilità di annullarla;
        \item Pubblicare e gestire recensioni su eventi a cui ha partecipato;
        \item Essere assegnato come organiser per un evento specifico;
    \end{itemize}
    \item \textbf{Organiser}: un utente registrato che assume il ruolo di organizzatore per un evento specifico, non rappresenta una tipologia 
    permanente di utente. L'organiser può, solamente per gli eventi a cui è stato assegnato, modificarne i dati, la programmazione 
    e le proprietà, oltre a eseguire le transizioni di stato permesse dal ciclo di vita dell'evento. L'organiser non può 
    creare eventi, nè assegnare o rimuovere altri organiser, tali sono operazioni riservate agli amministratori.
    \item \textbf{Amministratore}: un utente registrato con privilegi di amministrazione, che può:
    \begin{itemize}
        \item Creare e gestire venue e risorse;
        \item Creare eventi e risultarne il \textit{creator};
        \item Gestire qualsiasi evento;
        \item Assegnare e rimuovere organiser per qualsiasi evento;
        \item Prendere in carico, accettare o rifiutare richieste di evento;
        \item Moderare gli account ordinari degli utenti registrati;
        \item Confermare, cancellare o eliminare qualsiasi prenotazione
        \item Creare e consultare report di moderazione
    \end{itemize}
\end{itemize}
    
\textbf{Invited guest} è un caso particolare: si tratta di una persona inserita nella lista invitati di un evento privato. Essa non coincide necessariamente con 
un utente registrato, e non interagisce direttamente con il sistema, non si tratta dunque di un attore effettivo. La conferma o la cancellazione della propria partecipazione all'evento viene registrata 
direttamente dall'organiser o dall'amministratore, che possono anche modificare i dati anagrafici dell'invitato.



\section{Catalogo dei casi d'uso}
% !TEX root = ../main.tex
% Fonte: docs/USE-CASES.md, sezione 5.
\begingroup
\small
\setlength{\tabcolsep}{3pt}
\renewcommand{\arraystretch}{1.18}
\begin{longtable}{@{}>{\raggedright\arraybackslash}p{1.55cm}>{\raggedright\arraybackslash}p{5.15cm}>{\raggedright\arraybackslash}p{2.6cm}>{\raggedright\arraybackslash}p{1.65cm}>{\raggedright\arraybackslash}p{1.4cm}>{\raggedright\arraybackslash}p{2.3cm}@{}}
\caption{Catalogo dei casi d'uso.}\label{tab:catalogo-uc}\\
\toprule
\textbf{ID} & \textbf{Titolo} & \textbf{Attore principale} & \textbf{Livello} & \textbf{Priorità} & \textbf{Stato} \\
\midrule
\endfirsthead
\multicolumn{6}{l}{\small Catalogo dei casi d\textquotesingle uso (continua)}\\
\toprule
\textbf{ID} & \textbf{Titolo} & \textbf{Attore principale} & \textbf{Livello} & \textbf{Priorità} & \textbf{Stato} \\
\midrule
\endhead
\midrule
\multicolumn{6}{r}{\small Continua nella pagina successiva}\\
\endfoot
\bottomrule
\endlastfoot
UC-01 & Registrare un account & Visitor & User goal & Alta & Implementato \\
UC-02 & Gestire il proprio account & User & Summary & Alta & Implementato \\
UC-02.1 & Modificare il profilo & User & User goal & Alta & Implementato \\
UC-02.2 & Cambiare password & User & User goal & Alta & Implementato \\
UC-02.3 & Eliminare l'account & User & User goal & Media & Implementato \\
UC-03 & Moderare un account & Administrator & Summary & Alta & Implementato \\
UC-03.1 & Bannare un account & Administrator & User goal & Alta & Implementato \\
UC-03.2 & Riattivare un account & Administrator & User goal & Media & Implementato \\
UC-04 & Consultare venue, risorse ed eventi & User & Summary & Media & Implementato \\
UC-05 & Gestire una venue & Administrator & Summary & Alta & Implementato \\
UC-05.1 & Creare una venue & Administrator & User goal & Alta & Implementato \\
UC-05.2 & Modificare una venue & Administrator & User goal & Alta & Implementato \\
UC-05.3 & Eliminare una venue & Administrator & User goal & Media & Implementato \\
UC-06 & Gestire le risorse di una venue & Administrator & Summary & Alta & Implementato \\
UC-06.1 & Aggiungere una risorsa & Administrator & User goal & Alta & Implementato \\
UC-06.2 & Modificare una risorsa & Administrator & User goal & Alta & Implementato \\
UC-06.3 & Rimuovere una risorsa & Administrator & User goal & Media & Implementato \\
UC-07 & Presentare una richiesta di evento & User & User goal & Alta & Implementato \\
UC-08 & Gestire una richiesta di evento & Administrator & User goal & Alta & Implementato \\
UC-09 & Gestire un evento & Administrator / Organiser & Summary & Alta & Implementato \\
UC-09.1 & Creare un evento & Administrator & User goal & Alta & Implementato \\
UC-09.2 & Confermare un evento & Administrator / Organiser & User goal & Alta & Implementato \\
UC-09.3 & Pubblicare un evento & Administrator / Organiser & User goal & Alta & Implementato \\
UC-09.4 & Modificare o riprogrammare un evento & Administrator / Organiser & User goal & Media & Implementato \\
UC-09.5 & Cancellare un evento & Administrator / Organiser & User goal & Alta & Implementato \\
UC-09.6 & Assegnare o rimuovere un organiser & Administrator & User goal & Media & Implementato \\
UC-10 & Gestire la lista degli invitati & Administrator / Organiser & User goal & Media & Implementato \\
UC-11 & Prenotare biglietti & User & User goal & Alta & Implementato \\
UC-12 & Gestire una prenotazione & User / Administrator & Summary & Alta & Implementato \\
UC-12.1 & Confermare una prenotazione & User / Administrator & User goal & Alta & Implementato \\
UC-12.2 & Cancellare una prenotazione & User / Administrator & User goal & Alta & Implementato \\
UC-12.3 & Eliminare una prenotazione & User / Administrator & User goal & Media & Implementato \\
UC-13 & Pubblicare una recensione & User & User goal & Media & Implementato \\
UC-14 & Gestire le proprie recensioni & User & User goal & Media & Implementato \\
UC-15 & Creare e consultare report & Administrator & User goal & Media & Implementato \\
UC-16 & Autenticarsi & Visitor & Function & Alta & Implementato \\
\end{longtable}
\endgroup

UC-16 è una funzione di supporto: verifica l'identità dell'utente per consentire
l'accesso alle funzionalità applicative, senza concludere un obiettivo del dominio
come presentare una richiesta o ottenere una prenotazione.
Mantiene un identificatore e una rappresentazione autonoma nel diagramma del Visitor.
La versione corrente verifica credenziali e stato dell'account e restituisce lo User;
non crea sessioni né token.
\clearpage
\section{Template dei casi d'uso significativi}
% !TEX root = ../main.tex
% !TEX root = ../main.tex
% Fonte completa: docs/use-cases/UC-01-registrare-account.md
% Scheda compatta; per i test sono riportate le classi, non tutti i metodi.
\subsection{UC-01 - Registrare un account}
\label{uc:UC-01-registrare-account}
\begingroup
\fontsize{8.8}{10}\selectfont
\setlength{\tabcolsep}{5pt}
\setlength{\extrarowheight}{1pt}
\arrayrulecolor{lineeUC}
\setlength{\arrayrulewidth}{0.3pt}
\renewcommand{\arraystretch}{1.08}
\setlength{\parskip}{0pt}
\noindent\begin{tabularx}{\textwidth}{>{\columncolor{etichetteUC}\raggedright\arraybackslash\sffamily\bfseries\color{black!75}}p{2.5cm}>{\raggedright\arraybackslash}X}
\hline
ID e livello & UC-01 -- User goal \\[2pt] \hline
Attori & \textbf{Principale:} Visitor; \textbf{secondari:} Nessuno \\[2pt] \hline
Descrizione & Un Visitor richiede la registrazione, fornisce i propri dati e una password, e il sistema crea un nuovo account User ordinario.\par\textbf{Trigger:} Il Visitor richiede di registrare un nuovo account. \\[2pt] \hline
Precondizioni & \textbf{1.} Il Visitor non è autenticato. \enspace \textbf{2.} Il sistema è disponibile a ricevere la richiesta di registrazione. L'unicità dell'email e la validità dei dati non sono considerate precondizioni: vengono verificate durante il caso d'uso e, se non soddisfatte, producono un flusso alternativo. \\[2pt] \hline
Postcondizioni & \textbf{1.} È persistito un nuovo account User associato ai dati forniti. \enspace \textbf{2.} Il nuovo account non possiede privilegi da Administrator. \enspace \textbf{3.} Il nuovo account si trova nello stato \texttt{ACTIVE}. \enspace \textbf{4.} La password è persistita esclusivamente come hash PBKDF2 con salt casuale e non in chiaro. \enspace \textbf{5.} Il sistema restituisce l'identificativo del nuovo User. \\[2pt] \hline
Garanzia minima & Se la registrazione non può essere completata, non viene persistito alcun nuovo User e non viene memorizzata alcuna password. Eventuali valori amministrativi o di stato presenti nei dati ricevuti non permettono comunque di creare un Administrator o un account già \texttt{BANNED}. \\[2pt] \hline
Flusso principale & \begin{minipage}[t]{\linewidth}\begin{enumerate}[leftmargin=1.4em,nosep,label=\arabic*.]
\item Il Visitor richiede di registrare un nuovo account.
\item Il Visitor fornisce username, nome, cognome, data di nascita, email, eventuale numero di telefono e password.
\item Il sistema verifica la presenza e la validità formale dei dati forniti e della password.
\item Il sistema verifica che l'indirizzo email non sia già associato a un altro User.
\item Il sistema inizializza il nuovo account come User ordinario nello stato \texttt{ACTIVE}, ignorando eventuali valori amministrativi o di stato ricevuti.
\item Il sistema genera dalla password una rappresentazione sicura PBKDF2 con salt casuale.
\item Il sistema persiste atomicamente lo User e la credenziale codificata.
\item Il sistema conferma la transazione e restituisce l'identificativo del nuovo User.
\end{enumerate}\vspace{3pt}\end{minipage} \\[2pt] \hline
Flussi alternativi & \textbf{2a. Dati obbligatori mancanti} \par (1) Il sistema rileva che lo User da registrare o uno dei dati obbligatori non è stato fornito. (2) Il sistema rifiuta la richiesta come non valida. (3) Il caso d'uso termina senza persistere alcun User o password. \par\vspace{1pt} \textbf{3a. Dati personali non validi} \par (1) Il sistema rileva almeno una delle seguenti condizioni: \textbullet{} username, nome o cognome vuoti oppure con lunghezza non valida; \textbullet{} email con formato non valido; \textbullet{} numero di telefono facoltativo presente ma con formato o lunghezza non validi; \textbullet{} data di nascita futura, assente o precedente al 1900. (2) Il sistema rifiuta la richiesta indicando che i dati non sono validi. (3) Il caso d'uso termina senza persistere alcun User o password. \par\vspace{1pt} \textbf{3b. Password non valida} \par (1) Il sistema rileva che la password è nulla, vuota oppure non contiene tra 8 e 30 caratteri. (2) Il sistema rifiuta la richiesta. (3) Il caso d'uso termina senza persistere alcun User o password. \par\vspace{1pt} \textbf{4a. Email già utilizzata} \par (1) Il sistema rileva che l'indirizzo email è già associato a un altro User. (2) Il sistema rifiuta la registrazione per evitare la duplicazione dell'account. (3) Il caso d'uso termina senza persistere alcun nuovo User o password. \par\vspace{1pt} \textbf{7a. Conflitto o errore durante la persistenza} \par (1) Il sistema rileva che i dati non possono essere persistiti, ad esempio per un username già utilizzato o per un errore del database. (2) Il sistema annulla la transazione. (3) Nessun User incompleto e nessuna credenziale rimangono persistiti. (4) Il caso d'uso termina con un errore. \\[2pt] \hline
Regole & \textbullet{} \textbf{BR-01:} la registrazione crea sempre un utente ordinario \texttt{ACTIVE}; valori admin o banned forniti in input vengono ignorati. \textbullet{} \textbf{BR-02:} username, nome, cognome, email e data di nascita devono essere validi; l'email deve essere univoca. \textbullet{} \textbf{BR-03:} una password deve contenere tra 8 e 30 caratteri. \textbullet{} \textbf{BR-04:} le password vengono persistite esclusivamente come hash PBKDF2 con salt casuale. \\[2pt] \hline
Test & \texttt{User\allowbreak{}Service\allowbreak{}Workflow\allowbreak{}Test}; \texttt{User\allowbreak{}Service\allowbreak{}Validation\allowbreak{}Test}; \texttt{Password\allowbreak{}Hasher\allowbreak{}Test}; \texttt{Password\allowbreak{}Credential\allowbreak{}Integration\allowbreak{}Test}.\par Il vincolo di lunghezza della password, pur essendo implementato in \texttt{Auth\allowbreak{}Service}, non dispone ancora di un test unitario dedicato ai limiti di 8 e 30 caratteri. \\[2pt] \hline
\end{tabularx}\par
\endgroup

% !TEX root = ../main.tex
% Fonte completa: docs/use-cases/UC-03-1-bannare-account.md
% Scheda compatta; per i test sono riportate le classi, non tutti i metodi.
\clearpage
\subsection{UC-03.1 - Bannare un account}
\label{uc:UC-03-1-bannare-account}
\begingroup
\fontsize{9.2}{10.6}\selectfont
\setlength{\tabcolsep}{5pt}
\setlength{\extrarowheight}{2pt}
\arrayrulecolor{lineeUC}
\setlength{\arrayrulewidth}{0.3pt}
\renewcommand{\arraystretch}{1.12}
\setlength{\parskip}{0pt}
\noindent\begin{tabularx}{\textwidth}{>{\columncolor{etichetteUC}\raggedright\arraybackslash\sffamily\bfseries\color{black!75}}p{2.5cm}>{\raggedright\arraybackslash}X}
\hline
ID e livello & UC-03.1 -- User goal \\[2pt] \hline
Attori & \textbf{Principale:} Administrator; \textbf{secondari:} User \\[2pt] \hline
Descrizione & Un Administrator sospende un account ordinario attivo, impedendogli di utilizzare le operazioni non consentite agli utenti bannati.\par\textbf{Trigger:} L'Administrator richiede di bannare uno User. \\[2pt] \hline
Precondizioni & \textbf{1.} Esiste un account Administrator identificabile dal sistema. \par \textbf{2.} Lo User da moderare esiste. \\[2pt] \hline
Postcondizioni & \textbf{1.} Lo User si trova nello stato \texttt{BANNED}. \par \textbf{2.} Gli altri dati dello User rimangono invariati. \\[2pt] \hline
Garanzia minima & Se autenticazione, autorizzazione o transizione falliscono, lo stato dello User rimane invariato. \\[2pt] \hline
Flusso principale & \begin{minipage}[t]{\linewidth}\begin{enumerate}[leftmargin=1.4em,nosep,label=\arabic*.]
\item L'Administrator seleziona uno User attivo e richiede di bannarlo.
\item Il sistema verifica identificativo, password e ruolo dell'Administrator.
\item Il sistema carica lo User indicato.
\item Il sistema verifica che il bersaglio sia uno User ordinario diverso dall'attore.
\item Il sistema verifica che lo User sia \texttt{ACTIVE}.
\item Il sistema aggiorna atomicamente lo stato a \texttt{BANNED}.
\item Il sistema conferma l'operazione.
\end{enumerate}\vspace{3pt}\end{minipage} \\[2pt] \hline
Flussi alternativi & \textbf{2a. Credenziali amministrative non valide} \par (1) Il sistema nega l'operazione. (2) Il caso d'uso termina senza modifiche. \par\smallskip \textbf{3a. User inesistente} \par (1) Il sistema segnala che lo User non esiste. (2) Il caso d'uso termina senza modifiche. \par\smallskip \textbf{4a. Auto-moderazione o bersaglio Administrator} \par (1) Il sistema rileva che l'attore coincide con il bersaglio oppure che il bersaglio è Administrator. (2) Il sistema nega l'operazione. \par\smallskip \textbf{5a. User già bannato} \par (1) Il sistema rileva che lo User è già \texttt{BANNED}. (2) Il sistema rifiuta la transizione come conflitto. \\[2pt] \hline
Regole &\par\smallskip\textbullet{} \textbf{BR-05:} le operazioni amministrative richiedono identificativo e password validi di un Administrator.\par\smallskip\textbullet{} \textbf{BR-06:} un Administrator non può moderare se stesso né un altro Administrator. \\[2pt] \hline
Test & \texttt{User\allowbreak{}Service\allowbreak{}Workflow\allowbreak{}Test}. \\[2pt] \hline
\end{tabularx}\par
\endgroup

% !TEX root = ../main.tex
% Fonte completa: docs/use-cases/UC-03-2-riattivare-account.md
% Scheda compatta; per i test sono riportate le classi, non tutti i metodi.
\clearpage
\subsection{UC-03.2 - Riattivare un account}
\label{uc:UC-03-2-riattivare-account}
\begingroup
\fontsize{9.2}{10.6}\selectfont
\setlength{\tabcolsep}{5pt}
\setlength{\extrarowheight}{2pt}
\arrayrulecolor{lineeUC}
\setlength{\arrayrulewidth}{0.3pt}
\renewcommand{\arraystretch}{1.12}
\setlength{\parskip}{0pt}
\noindent\begin{tabularx}{\textwidth}{>{\columncolor{etichetteUC}\raggedright\arraybackslash\sffamily\bfseries\color{black!75}}p{2.5cm}>{\raggedright\arraybackslash}X}
\hline
ID e livello & UC-03.2 -- User goal \\[2pt] \hline
Attori & \textbf{Principale:} Administrator; \textbf{secondari:} User \\[2pt] \hline
Descrizione & Un Administrator riattiva un account ordinario precedentemente bannato.\par\textbf{Trigger:} L'Administrator richiede di riattivare uno User bannato. \\[2pt] \hline
Precondizioni & \textbf{1.} Esiste un account Administrator identificabile dal sistema. \par \textbf{2.} Lo User da moderare esiste. \\[2pt] \hline
Postcondizioni & \textbf{1.} Lo User si trova nello stato \texttt{ACTIVE}. \par \textbf{2.} Gli altri dati dello User rimangono invariati. \\[2pt] \hline
Garanzia minima & Se autenticazione, autorizzazione o transizione falliscono, lo stato dello User rimane invariato. \\[2pt] \hline
Flusso principale & \begin{minipage}[t]{\linewidth}\begin{enumerate}[leftmargin=1.4em,nosep,label=\arabic*.]
\item L'Administrator seleziona uno User bannato e richiede di riattivarlo.
\item Il sistema verifica identificativo, password e ruolo dell'Administrator.
\item Il sistema carica lo User indicato.
\item Il sistema verifica che il bersaglio sia uno User ordinario diverso dall'attore.
\item Il sistema verifica che lo User sia \texttt{BANNED}.
\item Il sistema aggiorna atomicamente lo stato a \texttt{ACTIVE}.
\item Il sistema conferma l'operazione.
\end{enumerate}\vspace{3pt}\end{minipage} \\[2pt] \hline
Flussi alternativi & \textbf{2a. Credenziali amministrative non valide} \par (1) Il sistema nega l'operazione. (2) Il caso d'uso termina senza modifiche. \par\smallskip \textbf{3a. User inesistente} \par (1) Il sistema segnala che lo User non esiste. (2) Il caso d'uso termina senza modifiche. \par\smallskip \textbf{4a. Auto-moderazione o bersaglio Administrator} \par (1) Il sistema rileva che l'attore coincide con il bersaglio oppure che il bersaglio è Administrator. (2) Il sistema nega l'operazione. \par\smallskip \textbf{5a. User già attivo} \par (1) Il sistema rileva che lo User è già \texttt{ACTIVE}. (2) Il sistema rifiuta la transizione come conflitto. \\[2pt] \hline
Regole &\par\smallskip\textbullet{} \textbf{BR-05:} le operazioni amministrative richiedono identificativo e password validi di un Administrator.\par\smallskip\textbullet{} \textbf{BR-06:} un Administrator non può moderare se stesso né un altro Administrator. \\[2pt] \hline
Test & \texttt{User\allowbreak{}Service\allowbreak{}Workflow\allowbreak{}Test}. \\[2pt] \hline
\end{tabularx}\par
\endgroup

% !TEX root = ../main.tex
% Fonte completa: docs/use-cases/UC-07-presentare-richiesta-evento.md
% Scheda compatta; per i test sono riportate le classi, non tutti i metodi.
\clearpage
\subsection{UC-07 - Presentare una richiesta di evento}
\label{uc:UC-07-presentare-richiesta-evento}
\begingroup
\fontsize{9.2}{10.6}\selectfont
\setlength{\tabcolsep}{5pt}
\setlength{\extrarowheight}{2pt}
\arrayrulecolor{lineeUC}
\setlength{\arrayrulewidth}{0.3pt}
\renewcommand{\arraystretch}{1.12}
\setlength{\parskip}{0pt}
\noindent\begin{tabularx}{\textwidth}{>{\columncolor{etichetteUC}\raggedright\arraybackslash\sffamily\bfseries\color{black!75}}p{2.5cm}>{\raggedright\arraybackslash}X}
\hline
ID e livello & UC-07 -- User goal \\[2pt] \hline
Attori & \textbf{Principale:} User; \textbf{secondari:} Nessuno \\[2pt] \hline
Descrizione & Uno User propone un evento indicando venue, nome, descrizione e intervallo temporale.\par\textbf{Trigger:} Lo User invia una nuova richiesta di evento. \\[2pt] \hline
Precondizioni & \textbf{1.} Lo User è registrato. \par \textbf{2.} La venue desiderata esiste. \\[2pt] \hline
Postcondizioni & \textbf{1.} La richiesta è persistita nello stato \texttt{PENDING}. \par \textbf{2.} La data di creazione è assegnata dal sistema. \par \textbf{3.} La richiesta non è ancora chiusa né necessariamente assegnata a un handler. \\[2pt] \hline
Garanzia minima & Una richiesta non valida non viene persistita. \\[2pt] \hline
Flusso principale & \begin{minipage}[t]{\linewidth}\begin{enumerate}[leftmargin=1.4em,nosep,label=\arabic*.]
\item Lo User seleziona una venue e inserisce nome, descrizione e date dell'evento proposto.
\item Il sistema verifica che il chiamante coincida con il richiedente, che questo esista e non sia Administrator.
\item Il sistema verifica l'esistenza della venue.
\item Il sistema valida nome, descrizione e intervallo temporale.
\item Il sistema verifica lo stato iniziale \texttt{PENDING} e, se assente, assegna l'istante corrente come data di creazione.
\item Il sistema persiste la richiesta e restituisce il suo identificativo.
\end{enumerate}\vspace{3pt}\end{minipage} \\[2pt] \hline
Flussi alternativi & \textbf{2a. Richiedente inesistente o non valido} \par (1) Il sistema rifiuta la richiesta. (2) Nessun dato viene persistito. \par\smallskip \textbf{2b. Richiedente Administrator} \par (1) Il sistema rileva che il richiedente non è uno User ordinario. (2) Il sistema rifiuta la richiesta. \par\smallskip \textbf{3a. Venue inesistente} \par (1) Il sistema segnala che la venue non esiste. (2) Il caso d'uso termina senza persistenza. \par\smallskip \textbf{4a. Dati descrittivi non validi} \par (1) Il sistema rileva un nome vuoto o non compreso tra 2 e 100 caratteri, oppure una descrizione superiore a 1000 caratteri. (2) Il sistema rifiuta la richiesta. \par\smallskip \textbf{4b. Intervallo temporale non valido} \par (1) Il sistema rileva una data mancante oppure un inizio non precedente alla fine. (2) Il sistema rifiuta la richiesta. \\[2pt] \hline
Regole &\par\smallskip\textbullet{} \textbf{BR-09:} il richiedente deve essere uno User ordinario, la venue deve esistere e nome e date devono essere validi.\par\smallskip\textbullet{} Una nuova richiesta entra nel ciclo di vita in stato \texttt{PENDING}. \\[2pt] \hline
Test & \texttt{Request\allowbreak{}And\allowbreak{}Guest\allowbreak{}Workflow\allowbreak{}Test}; \texttt{Secondary\allowbreak{}Service\allowbreak{}Validation\allowbreak{}Test}.\par Il sistema ignora eventuali handler, stato, preventivo e data di chiusura ricevuti in input, impedendo che la creazione aggiri il ciclo di vita della richiesta. \\[2pt] \hline
\end{tabularx}\par
\endgroup

% !TEX root = ../main.tex
% Fonte completa: docs/use-cases/UC-08-gestire-richiesta-evento.md
% Scheda compatta; per i test sono riportate le classi, non tutti i metodi.
\clearpage
\subsection{UC-08 - Gestire una richiesta di evento}
\label{uc:UC-08-gestire-richiesta-evento}
\begingroup
\fontsize{9.2}{10.6}\selectfont
\setlength{\tabcolsep}{5pt}
\setlength{\extrarowheight}{2pt}
\arrayrulecolor{lineeUC}
\setlength{\arrayrulewidth}{0.3pt}
\renewcommand{\arraystretch}{1.12}
\setlength{\parskip}{0pt}
\noindent\begin{tabularx}{\textwidth}{>{\columncolor{etichetteUC}\raggedright\arraybackslash\sffamily\bfseries\color{black!75}}p{2.5cm}>{\raggedright\arraybackslash}X}
\hline
ID e livello & UC-08 -- User goal \\[2pt] \hline
Attori & \textbf{Principale:} Administrator; \textbf{secondari:} User richiedente \\[2pt] \hline
Descrizione & Un Administrator prende in carico una richiesta pendente e la accetta o la rifiuta; il richiedente può cancellarla finché è pendente.\par\textbf{Trigger:} Un Administrator seleziona una richiesta \texttt{PENDING} da gestire. \\[2pt] \hline
Precondizioni & \textbf{1.} La richiesta esiste ed è \texttt{PENDING}. \par \textbf{2.} Per l'accettazione è stato assegnato un handler Administrator. \\[2pt] \hline
Postcondizioni & \textbf{1.} La richiesta è \texttt{ACCEPTED}, \texttt{REJECTED} oppure \texttt{CANCELLED}. \par \textbf{2.} Una richiesta accettata contiene un preventivo non negativo. \par \textbf{3.} La data di chiusura è assegnata dal sistema. \\[2pt] \hline
Garanzia minima & In caso di errore la richiesta conserva stato, handler, preventivo e data di chiusura precedenti. \\[2pt] \hline
Flusso principale & \begin{minipage}[t]{\linewidth}\begin{enumerate}[leftmargin=1.4em,nosep,label=\arabic*.]
\item L'Administrator seleziona una richiesta pendente.
\item Il sistema blocca la richiesta per evitare aggiornamenti concorrenti.
\item Il sistema verifica l'identità amministrativa del chiamante; l'Administrator assegna un handler Administrator, se non già presente.
\item L'Administrator sceglie di accettare la richiesta e indica il preventivo.
\item Il sistema verifica stato, handler e preventivo.
\item Il sistema imposta lo stato \texttt{ACCEPTED} e la data di chiusura corrente.
\item Il sistema persiste e conferma la transazione.
\end{enumerate}\vspace{3pt}\end{minipage} \\[2pt] \hline
Flussi alternativi & \textbf{3a. Handler inesistente o non Administrator} \par (1) Il sistema rifiuta l'assegnazione. (2) La richiesta rimane invariata. \par\smallskip \textbf{4a. Rifiuto della richiesta} \par (1) L'handler assegnato sceglie di rifiutare la richiesta pendente. (2) Il sistema imposta \texttt{REJECTED} e la data di chiusura corrente. (3) Il flusso riprende dal passo 7. \par\smallskip \textbf{4b. Cancellazione da parte del richiedente} \par (1) Il sistema verifica che il chiamante coincida con il richiedente, che cancella la propria richiesta ancora pendente. (2) Il sistema imposta \texttt{CANCELLED} e la data di chiusura corrente. (3) Il flusso riprende dal passo 7. \par\smallskip \textbf{5a. Preventivo non valido} \par (1) Il sistema rileva un preventivo assente o negativo. (2) L'accettazione viene rifiutata. \par\smallskip \textbf{5b. Richiesta già chiusa} \par (1) Il sistema rileva che la richiesta non è più \texttt{PENDING}. (2) La nuova transizione viene rifiutata come conflitto. \\[2pt] \hline
Regole &\par\smallskip\textbullet{} \textbf{BR-10:} soltanto uno User Administrator può essere assegnato come handler.\par\smallskip\textbullet{} \textbf{BR-11:} l'accettazione richiede stato \texttt{PENDING}, handler assegnato e preventivo non negativo.\par\smallskip\textbullet{} Le transizioni terminali sono \texttt{ACCEPTED}, \texttt{REJECTED} e \texttt{CANCELLED}. \\[2pt] \hline
Test & \texttt{Request\allowbreak{}And\allowbreak{}Guest\allowbreak{}Workflow\allowbreak{}Test}.\par Accettazione e rifiuto sono riservati all'handler assegnato, non a un Administrator qualsiasi. La cancellazione e le modifiche descrittive restano invece riservate al richiedente finché la richiesta è \texttt{PENDING}. \\[2pt] \hline
\end{tabularx}\par
\endgroup

% !TEX root = ../main.tex
% Fonte completa: docs/use-cases/UC-09-3-pubblicare-evento.md
% Scheda compatta; per i test sono riportate le classi, non tutti i metodi.
\clearpage
\subsection{UC-09.3 - Pubblicare un evento}
\label{uc:UC-09-3-pubblicare-evento}
\begingroup
\fontsize{9.2}{10.6}\selectfont
\setlength{\tabcolsep}{5pt}
\setlength{\extrarowheight}{2pt}
\arrayrulecolor{lineeUC}
\setlength{\arrayrulewidth}{0.3pt}
\renewcommand{\arraystretch}{1.12}
\setlength{\parskip}{0pt}
\noindent\begin{tabularx}{\textwidth}{>{\columncolor{etichetteUC}\raggedright\arraybackslash\sffamily\bfseries\color{black!75}}p{2.5cm}>{\raggedright\arraybackslash}X}
\hline
ID e livello & UC-09.3 -- User goal \\[2pt] \hline
Attori & \textbf{Principale:} Administrator / Organiser assegnato; \textbf{secondari:} Nessuno \\[2pt] \hline
Descrizione & Un attore autorizzato rende visibile e prenotabile un evento confermato e valido.\par\textbf{Trigger:} L'attore richiede la pubblicazione di un evento. \\[2pt] \hline
Precondizioni & \textbf{1.} L'evento esiste ed è \texttt{CONFIRMED}. \par \textbf{2.} L'attore è Administrator oppure Organiser assegnato all'evento. \\[2pt] \hline
Postcondizioni & \textbf{1.} L'evento è \texttt{PUBLISHED}. \par \textbf{2.} \texttt{published\allowbreak{}At} contiene l'istante della pubblicazione. \\[2pt] \hline
Garanzia minima & Se autorizzazione, stato o dati non sono validi, evento e data di pubblicazione rimangono invariati. \\[2pt] \hline
Flusso principale & \begin{minipage}[t]{\linewidth}\begin{enumerate}[leftmargin=1.4em,nosep,label=\arabic*.]
\item L'attore seleziona un evento confermato e ne richiede la pubblicazione.
\item Il sistema blocca l'evento e identifica l'attore.
\item Il sistema verifica che l'attore sia Administrator o Organiser assegnato.
\item Il sistema verifica la transizione \texttt{CONFIRMED} \ensuremath{\rightarrow} \texttt{PUBLISHED}.
\item Il sistema verifica che l'inizio sia futuro, la capienza positiva e la venue esistente.
\item Il sistema imposta stato e istante di pubblicazione.
\item Il sistema conferma la transazione.
\end{enumerate}\vspace{3pt}\end{minipage} \\[2pt] \hline
Flussi alternativi & \textbf{2a. Evento o attore inesistente} \par (1) Il sistema segnala l'elemento mancante. (2) Il caso d'uso termina senza modifiche. \par\smallskip \textbf{3a. Attore non autorizzato} \par (1) Il sistema rileva che lo User non è Administrator né Organiser assegnato. (2) La pubblicazione viene negata. \par\smallskip \textbf{4a. Stato incompatibile} \par (1) Il sistema rileva che l'evento non è \texttt{CONFIRMED}, oppure è già \texttt{PUBLISHED} o \texttt{CANCELLED}. (2) Il sistema rifiuta la transizione. \par\smallskip \textbf{5a. Evento non pubblicabile} \par (1) Il sistema rileva una data di inizio non futura, capienza non positiva o venue inesistente. (2) Il sistema rifiuta la pubblicazione. \\[2pt] \hline
Regole &\par\smallskip\textbullet{} \textbf{BR-16:} un evento è pubblicabile solo se confermato, futuro, con capienza positiva e venue esistente.\par\smallskip\textbullet{} L'Administrator può gestire ogni evento; l'Organiser soltanto quelli a lui assegnati. \\[2pt] \hline
Test & \texttt{Event\allowbreak{}Service\allowbreak{}Workflow\allowbreak{}Test}; \texttt{Status\allowbreak{}Transition\allowbreak{}Test}. \\[2pt] \hline
\end{tabularx}\par
\endgroup

% !TEX root = ../main.tex
% Fonte completa: docs/use-cases/UC-10-gestire-lista-invitati.md
% Scheda compatta; per i test sono riportate le classi, non tutti i metodi.
\clearpage
\subsection{UC-10 - Gestire la lista degli invitati}
\label{uc:UC-10-gestire-lista-invitati}
\begingroup
\fontsize{8.8}{10.0}\selectfont
\setlength{\tabcolsep}{5pt}
\setlength{\extrarowheight}{1pt}
\arrayrulecolor{lineeUC}
\setlength{\arrayrulewidth}{0.3pt}
\renewcommand{\arraystretch}{1.12}
\setlength{\parskip}{0pt}
\noindent\begin{tabularx}{\textwidth}{>{\columncolor{etichetteUC}\raggedright\arraybackslash\sffamily\bfseries\color{black!75}}p{2.5cm}>{\raggedright\arraybackslash}X}
\hline
ID e livello & UC-10 -- User goal \\[2pt] \hline
Attori & \textbf{Principale:} Administrator / Organiser assegnato; \textbf{secondari:} Invited Guest \\[2pt] \hline
Descrizione & Un attore autorizzato inserisce, aggiorna, conferma, cancella o rimuove invitati da un evento privato.\par\textbf{Trigger:} L'attore apre la lista invitati di un evento. \\[2pt] \hline
Precondizioni & \textbf{1.} L'evento esiste, non è \texttt{CANCELLED} ed è \texttt{PRIVATE\_\allowbreak{}GUEST\_\allowbreak{}LIST}. \par \textbf{2.} L'attore è Administrator oppure Organiser assegnato. \\[2pt] \hline
Postcondizioni & \textbf{1.} La lista riflette l'operazione richiesta. \par \textbf{2.} Un nuovo invitato nasce sempre \texttt{INVITED}. \par \textbf{3.} Evento e stato non possono essere alterati indirettamente tramite l'aggiornamento anagrafico. \\[2pt] \hline
Garanzia minima & Una richiesta non autorizzata o una transizione non valida non modifica la lista. \\[2pt] \hline
Flusso principale & \begin{minipage}[t]{\linewidth}\begin{enumerate}[leftmargin=1.4em,nosep,label=\arabic*.]
\item L'attore seleziona un evento privato.
\item Il sistema blocca l'evento e verifica l'autorizzazione.
\item L'attore fornisce nome e cognome del nuovo invitato.
\item Il sistema valida i dati e forza lo stato \texttt{INVITED}.
\item Il sistema persiste l'invitato.
\item Il sistema conferma l'operazione.
\end{enumerate}\vspace{3pt}\end{minipage} \\[2pt] \hline
Flussi alternativi & \textbf{2a. Evento non idoneo} \par (1) Il sistema rileva un evento inesistente, pubblico, con altra visibilità oppure cancellato. (2) L'operazione viene rifiutata. \par\smallskip \textbf{2b. Attore non autorizzato} \par (1) Il sistema rileva che l'attore non è Administrator né Organiser assegnato. (2) L'operazione viene negata. \par\smallskip \textbf{3a. Aggiornamento dei dati dell'invitato} \par (1) L'attore modifica i dati anagrafici di un invitato esistente. (2) Il sistema conserva evento e stato memorizzati e aggiorna soltanto i dati ammessi. \par\smallskip \textbf{3b. Conferma o cancellazione dell'invito} \par (1) L'attore richiede una transizione di stato. (2) Il sistema blocca invitato ed evento e verifica la transizione. (3) Sono consentite \texttt{INVITED} \ensuremath{\rightarrow} \texttt{CONFIRMED}, \texttt{INVITED} \ensuremath{\rightarrow} \texttt{CANCELLED} e \texttt{CONFIRMED} \ensuremath{\rightarrow} \texttt{CANCELLED}. (4) Il sistema persiste il nuovo stato. \par\smallskip \textbf{3c. Rimozione dell'invitato} \par (1) L'attore richiede la rimozione. (2) Il sistema verifica nuovamente autorizzazione ed evento, quindi elimina l'invitato. \par\smallskip \textbf{3d. Dati o transizione non validi} \par (1) Il sistema rileva dati obbligatori mancanti, una duplicazione di stato o una riattivazione da \texttt{CANCELLED}. (2) L'operazione viene rifiutata senza modifiche. \\[2pt] \hline
Regole &\par\smallskip\textbullet{} \textbf{BR-25:} gli invitati appartengono soltanto a eventi \texttt{PRIVATE\_\allowbreak{}GUEST\_\allowbreak{}LIST} e iniziano in stato \texttt{INVITED}.\par\smallskip\textbullet{} L'Organiser può gestire esclusivamente gli eventi a lui assegnati; l'Administrator può gestirli tutti. \\[2pt] \hline
Test & \texttt{Request\allowbreak{}And\allowbreak{}Guest\allowbreak{}Workflow\allowbreak{}Test}; \texttt{Event\allowbreak{}Guest\allowbreak{}Service\allowbreak{}Validation\allowbreak{}Test}; \texttt{Status\allowbreak{}Transition\allowbreak{}Test}. \\[2pt] \hline
Osservazioni & L'Invited Guest è un soggetto interessato, ma non un attore software: non usa direttamente il sistema. Administrator o Organiser registrano il suo riscontro. \\[2pt] \hline
\end{tabularx}\par
\endgroup

% !TEX root = ../main.tex
% Fonte completa: docs/use-cases/UC-11-prenotare-biglietti.md
% Scheda compatta; per i test sono riportate le classi, non tutti i metodi.
\clearpage
\subsection{UC-11 - Prenotare biglietti}
\label{uc:UC-11-prenotare-biglietti}
\begingroup
\fontsize{8.3}{9.1}\selectfont
\setlength{\tabcolsep}{5pt}
\setlength{\extrarowheight}{1pt}
\arrayrulecolor{lineeUC}
\setlength{\arrayrulewidth}{0.3pt}
\renewcommand{\arraystretch}{1}
\setlength{\parskip}{0pt}
\noindent\begin{tabularx}{\textwidth}{>{\columncolor{etichetteUC}\raggedright\arraybackslash\sffamily\bfseries\color{black!75}}p{2.5cm}>{\raggedright\arraybackslash}X}
\hline
ID e livello & UC-11 -- User goal \\[1pt] \hline
Attori & \textbf{Principale:} User; \textbf{secondari:} Nessuno \\[1pt] \hline
Descrizione & Uno User richiede uno o più biglietti nominativi per un evento pubblico, pubblicato e non ancora iniziato. Il sistema verifica utente, evento, dati dei biglietti e capienza residua, quindi crea atomicamente la prenotazione e i relativi biglietti.\par\textbf{Trigger:} Lo User richiede di prenotare uno o più biglietti per un evento selezionato. \\[1pt] \hline
Precondizioni & \textbf{1.} Lo User dispone di un account registrato e identificabile dal sistema. \enspace \textbf{2.} L'evento selezionato esiste. \enspace \textbf{3.} Lo User fornisce almeno un biglietto con nome e cognome dell'intestatario. Le condizioni relative allo stato dell'account, alla prenotabilità dell'evento e alla capienza vengono nuovamente verificate dal sistema durante l'operazione, perché possono cambiare nel tempo. \\[1pt] \hline
Postcondizioni & \textbf{1.} È persistita una nuova prenotazione associata allo User e all'evento. \enspace \textbf{2.} La prenotazione si trova nello stato \texttt{PENDING\_\allowbreak{}PAYMENT}. \enspace \textbf{3.} Il prezzo totale equivale al numero di biglietti moltiplicato per il prezzo unitario dell'evento; per un evento gratuito il totale è zero. \enspace \textbf{4.} Tutti i biglietti richiesti sono persistiti e associati alla nuova prenotazione. \enspace \textbf{5.} Ogni biglietto riporta come orario di inizio quello dell'evento. \\[1pt] \hline
Garanzia minima & Se l'operazione non può essere completata, non viene persistita alcuna prenotazione parziale e non viene inserito alcun biglietto. Tutte le scritture appartengono alla stessa transazione. \\[1pt] \hline
Flusso principale & \begin{minipage}[t]{\linewidth}\begin{enumerate}[leftmargin=1.4em,nosep,label=\arabic*.]
\item Lo User seleziona un evento e richiede uno o più biglietti nominativi.
\item Il sistema verifica che la lista dei biglietti non sia nulla o vuota e che ogni biglietto contenga nome e cognome validi.
\item Il sistema verifica che lo User esista e che il suo account non sia \texttt{BANNED}.
\item Il sistema carica e blocca l'evento per impedire che prenotazioni concorrenti modifichino contemporaneamente la disponibilità considerata.
\item Il sistema verifica che l'evento sia \texttt{PUBLISHED}, abbia visibilità \texttt{PUBLIC} e non sia ancora iniziato.
\item Il sistema conta i biglietti già emessi per l'evento.
\item Il sistema verifica che la capienza residua sia sufficiente per tutti i biglietti richiesti.
\item Il sistema calcola il prezzo totale della prenotazione.
\item Il sistema crea la prenotazione nello stato \texttt{PENDING\_\allowbreak{}PAYMENT} e la persiste.
\item Il sistema associa ogni biglietto alla nuova prenotazione e gli assegna l'orario iniziale dell'evento.
\item Il sistema persiste tutti i biglietti richiesti.
\item Il sistema conferma la transazione e restituisce la prenotazione creata.
\end{enumerate}\vspace{3pt}\end{minipage} \\[1pt] \hline
Flussi alternativi & \textbf{2a. Lista dei biglietti assente o vuota} \par (1) Il sistema rifiuta la richiesta come non valida. (2) Il caso d'uso termina senza modificare i dati persistiti. \par\vspace{1pt} \textbf{2b. Dati di un biglietto non validi} \par (1) Il sistema rileva un biglietto nullo, già associato a una prenotazione oppure privo di nome o cognome validi. (2) Il sistema rifiuta l'intera richiesta. (3) Il caso d'uso termina senza modificare i dati persistiti. \par\vspace{1pt} \textbf{3a. User inesistente} \par (1) Il sistema segnala che lo User non esiste. (2) L'evento non viene bloccato e il caso d'uso termina. \par\vspace{1pt} \textbf{3b. User bannato} \par (1) Il sistema nega l'operazione. (2) L'evento non viene bloccato e il caso d'uso termina. \par\vspace{1pt} \textbf{4a. Evento inesistente} \par (1) Il sistema segnala che l'evento non esiste. (2) Il caso d'uso termina senza modificare i dati persistiti. \par\vspace{1pt} \textbf{5a. Evento non prenotabile} \par (1) Il sistema rileva che l'evento non è \texttt{PUBLISHED}, non è \texttt{PUBLIC} oppure è già iniziato. (2) Il sistema rifiuta la richiesta. (3) Il caso d'uso termina senza modificare i dati persistiti. \par\vspace{1pt} \textbf{7a. Capienza insufficiente} \par (1) Il sistema rileva che i posti residui non sono sufficienti per tutti i biglietti richiesti. (2) Il sistema rifiuta l'intera richiesta: non è prevista una prenotazione parziale. (3) Il caso d'uso termina senza inserire booking o ticket. \par\vspace{1pt} \textbf{9a-11a. Errore durante la persistenza} \par (1) Il sistema annulla la transazione. (2) L'eventuale prenotazione e gli eventuali biglietti inseriti durante il tentativo non restano persistiti. (3) Il caso d'uso termina con un errore. \\[1pt] \hline
Regole & \textbullet{} \textbf{BR-07:} un utente \texttt{BANNED} non può prenotare eventi. \textbullet{} \textbf{BR-19:} sono prenotabili soltanto eventi \texttt{PUBLISHED}, \texttt{PUBLIC} e non ancora iniziati. \textbullet{} \textbf{BR-20:} una prenotazione deve contenere almeno un biglietto valido e nominativo. \textbullet{} \textbf{BR-21:} il prezzo totale è \texttt{numero biglietti \ensuremath{\times} prezzo unitario}; un prezzo nullo equivale a zero. \textbullet{} \textbf{BR-22:} il controllo della capienza e la scrittura di booking e ticket avvengono nella stessa transazione. \textbullet{} \textbf{BR-23:} l'evento viene letto con lock prima del conteggio dei biglietti per impedire overbooking concorrente. \textbullet{} \textbf{BR-24:} ogni ticket viene associato alla prenotazione appena creata e riceve l'orario iniziale dell'evento. \\[1pt] \hline
Test & \texttt{Booking\allowbreak{}Service\allowbreak{}Validation\allowbreak{}Test}; \texttt{Core\allowbreak{}Service\allowbreak{}Workflow\allowbreak{}Test}. \\[1pt] \hline
Osservazioni & Lo stato \texttt{PENDING\_\allowbreak{}PAYMENT} rappresenta soltanto una fase del ciclo applicativo della prenotazione. Il sistema non integra un gateway di pagamento: la successiva conferma appartiene a UC-12.1 e non al presente caso d'uso. \\[1pt] \hline
\end{tabularx}\par
\endgroup

% !TEX root = ../main.tex
% Fonte completa: docs/use-cases/UC-13-pubblicare-recensione.md
% Scheda compatta; per i test sono riportate le classi, non tutti i metodi.
\clearpage
\subsection{UC-13 - Pubblicare una recensione}
\label{uc:UC-13-pubblicare-recensione}
\begingroup
\fontsize{9.2}{10.6}\selectfont
\setlength{\tabcolsep}{5pt}
\setlength{\extrarowheight}{2pt}
\arrayrulecolor{lineeUC}
\setlength{\arrayrulewidth}{0.3pt}
\renewcommand{\arraystretch}{1.12}
\setlength{\parskip}{0pt}
\noindent\begin{tabularx}{\textwidth}{>{\columncolor{etichetteUC}\raggedright\arraybackslash\sffamily\bfseries\color{black!75}}p{2.5cm}>{\raggedright\arraybackslash}X}
\hline
ID e livello & UC-13 -- User goal \\[2pt] \hline
Attori & \textbf{Principale:} User; \textbf{secondari:} Nessuno \\[2pt] \hline
Descrizione & Uno User che ha partecipato a un evento concluso pubblica una sola recensione con voto ed eventuale commento.\par\textbf{Trigger:} Lo User richiede di recensire un evento concluso. \\[2pt] \hline
Precondizioni & \textbf{1.} Lo User possiede una prenotazione \texttt{CONFIRMED} per l'evento. \par \textbf{2.} L'evento esiste ed è terminato. \\[2pt] \hline
Postcondizioni & \textbf{1.} È persistita una recensione associata allo User e all'evento. \par \textbf{2.} La data di creazione è assegnata dal sistema. \\[2pt] \hline
Garanzia minima & Se un controllo fallisce non viene inserita alcuna recensione. \\[2pt] \hline
Flusso principale & \begin{minipage}[t]{\linewidth}\begin{enumerate}[leftmargin=1.4em,nosep,label=\arabic*.]
\item Lo User seleziona un evento terminato e fornisce voto ed eventuale commento.
\item Il sistema verifica che l'attore coincida con l'autore indicato.
\item Il sistema valida voto e commento.
\item Il sistema verifica l'esistenza e la conclusione dell'evento.
\item Il sistema verifica una prenotazione \texttt{CONFIRMED} dello User per l'evento.
\item Il sistema verifica che lo User non abbia già recensito l'evento.
\item Il sistema assegna la data corrente e persiste la recensione.
\end{enumerate}\vspace{3pt}\end{minipage} \\[2pt] \hline
Flussi alternativi & \textbf{2a. Attore diverso dall'autore} \par (1) Il sistema nega l'operazione. (2) Nessuna recensione viene inserita. \par\smallskip \textbf{3a. Voto o commento non validi} \par (1) Il sistema rileva un voto fuori dall'intervallo 1-5 o un commento superiore a 1000 caratteri. (2) La richiesta viene rifiutata. \par\smallskip \textbf{4a. Evento inesistente o non concluso} \par (1) Il sistema segnala l'evento mancante oppure non ancora terminato. (2) Il caso d'uso termina senza persistenza. \par\smallskip \textbf{5a. Partecipazione non dimostrata} \par (1) Il sistema non trova una prenotazione \texttt{CONFIRMED} per lo User e l'evento. (2) La pubblicazione viene negata. \par\smallskip \textbf{6a. Recensione già presente} \par (1) Il sistema rileva una recensione dello stesso User per lo stesso evento. (2) Il sistema rifiuta il duplicato come conflitto. \\[2pt] \hline
Regole &\par\smallskip\textbullet{} \textbf{BR-26:} il voto deve essere compreso tra 1 e 5 e il commento non può superare 1000 caratteri.\par\smallskip\textbullet{} \textbf{BR-27:} può recensire soltanto chi possiede una prenotazione \texttt{CONFIRMED} per un evento concluso.\par\smallskip\textbullet{} Uno User può pubblicare una sola recensione per evento e soltanto a proprio nome. \\[2pt] \hline
Test & \texttt{Secondary\allowbreak{}Service\allowbreak{}Workflow\allowbreak{}Test}; \texttt{Secondary\allowbreak{}Service\allowbreak{}Validation\allowbreak{}Test}; \texttt{Postgres\allowbreak{}Review\allowbreak{}Constraint\allowbreak{}Integration\allowbreak{}Test}. \\[2pt] \hline
\end{tabularx}\par
\endgroup

% !TEX root = ../main.tex
% Fonte completa: docs/use-cases/UC-15-creare-consultare-report.md
% Scheda compatta; per i test sono riportate le classi, non tutti i metodi.
\clearpage
\subsection{UC-15 - Creare e consultare report}
\label{uc:UC-15-creare-consultare-report}
\begingroup
\fontsize{9.2}{10.6}\selectfont
\setlength{\tabcolsep}{5pt}
\setlength{\extrarowheight}{2pt}
\arrayrulecolor{lineeUC}
\setlength{\arrayrulewidth}{0.3pt}
\renewcommand{\arraystretch}{1.12}
\setlength{\parskip}{0pt}
\noindent\begin{tabularx}{\textwidth}{>{\columncolor{etichetteUC}\raggedright\arraybackslash\sffamily\bfseries\color{black!75}}p{2.5cm}>{\raggedright\arraybackslash}X}
\hline
ID e livello & UC-15 -- User goal \\[2pt] \hline
Attori & \textbf{Principale:} Administrator; \textbf{secondari:} User segnalato \\[2pt] \hline
Descrizione & Un Administrator registra un report relativo a uno User ordinario, eventualmente collegato a un evento, e consulta i report tramite diversi filtri.\par\textbf{Trigger:} L'Administrator richiede di registrare o consultare un report. \\[2pt] \hline
Precondizioni & \textbf{1.} L'Administrator indicato esiste e possiede privilegi amministrativi. \par \textbf{2.} Lo User segnalato esiste e non è Administrator. \\[2pt] \hline
Postcondizioni & \textbf{1.} In creazione, il report è persistito con severità, commento, riferimenti e data assegnata dal sistema. \par \textbf{2.} In consultazione, il sistema restituisce i report corrispondenti al filtro senza modificare dati. \\[2pt] \hline
Garanzia minima & Un report non valido non viene persistito; una consultazione non produce scritture. \\[2pt] \hline
Flusso principale & \begin{minipage}[t]{\linewidth}\begin{enumerate}[leftmargin=1.4em,nosep,label=\arabic*.]
\item L'Administrator seleziona lo User da segnalare.
\item Indica severità, eventuale commento ed eventuale evento correlato.
\item Il sistema verifica che il chiamante sia Administrator e coincida con l'autore indicato.
\item Il sistema verifica che il bersaglio sia uno User ordinario.
\item Se presente, il sistema verifica l'esistenza dell'evento.
\item Il sistema valida severità e commento.
\item Il sistema assegna la data corrente e persiste il report.
\item Il sistema restituisce l'identificativo generato.
\end{enumerate}\vspace{3pt}\end{minipage} \\[2pt] \hline
Flussi alternativi & \textbf{1a. Consultazione dei report} \par (1) L'Administrator sceglie un criterio: identificativo, User, Administrator autore, evento, severità o loro combinazione supportata. (2) Il sistema verifica che il chiamante sia Administrator ed esegue una transazione read-only. (3) Il sistema restituisce i risultati corrispondenti. \par\smallskip \textbf{3a. Autore inesistente o non Administrator} \par (1) Il sistema segnala l'autore mancante oppure nega l'operazione per ruolo insufficiente. (2) Nessun report viene inserito. \par\smallskip \textbf{4a. Bersaglio inesistente o Administrator} \par (1) Il sistema segnala lo User mancante oppure rifiuta la segnalazione di un Administrator. (2) Nessun report viene inserito. \par\smallskip \textbf{5a. Evento correlato inesistente} \par (1) Il sistema segnala che l'evento indicato non esiste. (2) Il caso d'uso termina senza persistenza. \par\smallskip \textbf{6a. Contenuto non valido} \par (1) Il sistema rileva severità assente oppure commento superiore a 1000 caratteri. (2) La richiesta viene rifiutata. \\[2pt] \hline
Regole &\par\smallskip\textbullet{} \textbf{BR-29:} un report è creato da un Administrator, riguarda uno User ordinario e può riferirsi facoltativamente a un evento esistente.\par\smallskip\textbullet{} La severità è obbligatoria; il commento facoltativo non può superare 1000 caratteri. \\[2pt] \hline
Test & \texttt{Secondary\allowbreak{}Service\allowbreak{}Workflow\allowbreak{}Test}; \texttt{Secondary\allowbreak{}Service\allowbreak{}Validation\allowbreak{}Test}; \texttt{Service\allowbreak{}Query\allowbreak{}Delegation\allowbreak{}Test}.\par Creazione, consultazione, aggiornamento ed eliminazione ricevono separatamente l'identità del chiamante e sono riservate agli Administrator. In creazione, il chiamante deve inoltre coincidere con l'autore registrato nel report. \\[2pt] \hline
\end{tabularx}\par
\endgroup


\clearpage
\section{Regole di business e requisiti di qualità}
% !TEX root = ../main.tex
% Fonte: docs/SPECIFICHE_CONDIVISE.md, sezione 8.
\begingroup
\small
\setlength{\tabcolsep}{5pt}
\renewcommand{\arraystretch}{1.12}
\begin{longtable}{@{}>{\raggedright\arraybackslash\bfseries}p{1.4cm}>{\raggedright\arraybackslash}p{\dimexpr\textwidth-1.4cm-2\tabcolsep\relax}@{}}
\caption{Regole di business.}\label{tab:business-rules}\\
\toprule
\textbf{ID} & \textbf{Regola} \\
\midrule
\endfirsthead
\multicolumn{2}{@{}l}{\small Regole di business (continua)}\\
\toprule
\textbf{ID} & \textbf{Regola} \\
\midrule
\endhead
\midrule
\multicolumn{2}{r@{}}{\small Continua nella pagina successiva}\\
\endfoot
\bottomrule
\endlastfoot
\rowcolor{etichetteUC}
\multicolumn{2}{@{}l}{\rule{0pt}{2.6ex}\textbf{Utenti e sicurezza}} \\*[3pt]
\strut\label{br:BR-01}BR-01 & la registrazione crea sempre un utente ordinario \texttt{ACTIVE}; valori admin o banned forniti in input vengono ignorati. \\[4pt]
\strut\label{br:BR-02}BR-02 & username, nome, cognome, email e data di nascita devono essere validi; l'email deve essere univoca. \\[4pt]
\strut\label{br:BR-03}BR-03 & una password deve contenere tra 8 e 30 caratteri. \\[4pt]
\strut\label{br:BR-04}BR-04 & le password vengono persistite esclusivamente come hash PBKDF2 con salt casuale. \\[4pt]
\strut\label{br:BR-05}BR-05 & autenticazione, modifica profilo, eliminazione account, cambio password e moderazione richiedono la verifica delle credenziali appropriate; gli account \texttt{BANNED} non possono autenticarsi. \\[4pt]
\strut\label{br:BR-06}BR-06 & un admin non può bannare sé stesso o un altro admin. \\[4pt]
\strut\label{br:BR-07}BR-07 & un utente \texttt{BANNED} non può prenotare eventi. \\[4pt]
\rowcolor{etichetteUC}
\multicolumn{2}{@{}l}{\rule{0pt}{2.6ex}\textbf{Venue, richieste ed eventi}} \\*[3pt]
\strut\label{br:BR-08}BR-08 & una venue richiede nome e indirizzo strutturato completo; \texttt{additional\_info} è facoltativo. \\[4pt]
\strut\label{br:BR-09}BR-09 & una richiesta richiede requester ordinario, venue esistente, nome e intervallo temporale valido. \\[4pt]
\strut\label{br:BR-10}BR-10 & soltanto un admin può essere assegnato come handler. \\[4pt]
\strut\label{br:BR-11}BR-11 & una richiesta può essere accettata solo se ancora \texttt{PENDING}, con handler assegnato e preventivo non negativo. \\[4pt]
\strut\label{br:BR-12}BR-12 & un evento nuovo viene forzato a \texttt{DRAFT}, indipendentemente dallo stato ricevuto in input. \\[4pt]
\strut\label{br:BR-13}BR-13 & un evento richiede venue e creator esistenti, data finale successiva a quella iniziale e capienza positiva. \\[4pt]
\strut\label{br:BR-14}BR-14 & il creator deve essere un admin; l'organiser, quando presente, deve essere un utente ordinario. \\[4pt]
\strut\label{br:BR-15}BR-15 & organiser e poster sono facoltativi; il prezzo può essere nullo, interpretato come evento gratuito. \\[4pt]
\strut\label{br:BR-16}BR-16 & per pubblicare sono necessari stato \texttt{CONFIRMED}, dati validi, venue valida, capienza valida e coerenza temporale. \\[4pt]
\strut\label{br:BR-17}BR-17 & la capienza non può essere ridotta sotto il numero di biglietti già emessi. \\[4pt]
\strut\label{br:BR-18}BR-18 & cancellare un evento annulla nella stessa transazione prenotazioni e inviti attivi. \\[4pt]
\rowcolor{etichetteUC}
\multicolumn{2}{@{}l}{\rule{0pt}{2.6ex}\textbf{Prenotazioni e invitati}} \\*[3pt]
\strut\label{br:BR-19}BR-19 & sono prenotabili soltanto eventi \texttt{PUBLISHED}, \texttt{PUBLIC} e non ancora iniziati. \\[4pt]
\strut\label{br:BR-20}BR-20 & una prenotazione deve contenere almeno un biglietto valido e nominativo. \\[4pt]
\strut\label{br:BR-21}BR-21 & il prezzo totale è \texttt{numero biglietti \ensuremath{\times} prezzo unitario}; un prezzo nullo equivale a zero. \\[4pt]
\strut\label{br:BR-22}BR-22 & il controllo della capienza e la scrittura di booking e ticket avvengono nella stessa transazione. \\[4pt]
\strut\label{br:BR-23}BR-23 & l'evento viene letto con lock prima del conteggio dei biglietti per impedire overbooking concorrente. \\[4pt]
\strut\label{br:BR-24}BR-24 & ogni ticket viene associato alla prenotazione appena creata e riceve l'orario iniziale dell'evento. \\[4pt]
\strut\label{br:BR-25}BR-25 & gli invitati appartengono soltanto a eventi privati con guest list; il loro stato iniziale è sempre \texttt{INVITED}. \\[4pt]
\rowcolor{etichetteUC}
\multicolumn{2}{@{}l}{\rule{0pt}{2.6ex}\textbf{Feedback e moderazione}} \\*[3pt]
\strut\label{br:BR-26}BR-26 & una recensione ha rating compreso tra 1 e 5 e commento facoltativo fino a 1000 caratteri. \\[4pt]
\strut\label{br:BR-27}BR-27 & si può recensire soltanto un evento concluso per cui l'utente possiede almeno una prenotazione \texttt{CONFIRMED}. \\[4pt]
\strut\label{br:BR-28}BR-28 & un utente può recensire lo stesso evento una sola volta; il vincolo esiste sia nel service sia nel database. \\[4pt]
\strut\label{br:BR-29}BR-29 & attualmente soltanto un admin può creare un report e il soggetto segnalato deve essere un utente non admin. \\[4pt]
\strut\label{br:BR-30}BR-30 & un report può riferirsi facoltativamente a un evento e possiede severità \texttt{LOW}, \texttt{MIDDLE} o \texttt{HIGH}. \\[4pt]
\end{longtable}
\endgroup

